\section{Controllers}
\label{sec:controllers}

This section describes the three parts of the controller the project
designed: the flight layer that lets a drone follow a flocking model, the three
models, and the leaders' planning layer. Each part reached its form through
failures observed in simulation; the main ones are told here, because they
explain the design.

\subsection{Flight: the ideal boid is not the robot}
\label{sec:flight}

The first version turned each boid acceleration straight into a tilt and
flew Reynolds' rules unchanged. Twenty drones hovered calmly, then
neighbouring pairs started swinging in and out with a period of about 2\,s,
wider each time, until two touched and the physics blew up. The rules assume
a boid accelerates the moment it is asked to; a quadrotor has to tilt first,
and the attitude controller, with the gains of the Webots example, was slow
and underdamped (a tilt step overshot by 45\%). A stiff spring, separation
growing as $1/d^2$, behind a lagging, resonant actuator oscillates.
\citet{viragh2014} report the same delay-induced oscillations in flying
robots, colliding even in a hovering flock, and describe a robot by exactly
such imperfections.

So the drone was made to follow first, before touching the rules:
\begin{itemize}
  \item attitude is held by two nested loops, as in the Crazyflie firmware:
        the angle error sets a wanted rotation rate, capped at 1.5\,rad/s,
        and the rate error sets the motor differential. Faster rotations
        flipped the drone in hard reversals. A 21\textdegree\ tilt step now
        reaches 90\% in 0.26\,s with 4\% overshoot, where before it never
        did within a second;
  \item the hover speed comes from the physics ($4k\omega^2 = mg$, so
        $\omega = 55.4$\,rad/s, against 55.37 measured), and the motors speed
        up by $1/\sqrt{\cos\phi\cos\theta}$ when tilted;
  \item every model outputs a wanted velocity $\mathbf v_d$, which a velocity
        layer flies as $\mathbf a = (\mathbf v_d - \mathbf v)/\tau$, capped at
        $a_{\max} = 2.5$\,m/s$^2$. Velocity steps measured against a point
        mass with the same $\tau$ showed that below $\tau \approx 0.3$\,s the
        response stops getting faster, since that is the attitude loops' own
        speed, and only overshoots; $\tau = 0.3$\,s is the fastest setting
        without overshoot.
\end{itemize}
The attitude loops alone were enough for the same rules, unchanged, to settle
into a calm lattice with no collision (closest pair 0.64\,m, against contact
before). The velocity layer came later, with the hybrid model, whose terms are
velocities; it describes the robot by two measured numbers, $\tau$ and
$a_{\max}$, so every model meets the same robot.

\subsection{Three flocking models}
\label{sec:models}

\paragraph{Notation.} For drone $i$, $\mathcal N_i$ are the neighbours it
hears within the perception radius $\rho = 2.5$\,m; $\mathbf d_{ij}$ is the
offset from $i$ to $j$, $d_{ij} = |\mathbf d_{ij}|$, and $\hat{\mathbf x}$ is
the unit vector of $\mathbf x$. $\mathbf v_i$ is a drone's velocity and
$\mathbf g_i$, for a leader, the direction to go, from the route
(\cref{sec:planner}) or straight at C, with $|\mathbf g_i|$ the distance left.
$v_0 = 1$\,m/s is the cruise speed. Every speed in the models is a multiple of
$v_0$ (rates scale with it, decelerations with its square), so the rules are
exactly speed-free: at $k$ times $v_0$ the same flock flies the same paths, $k$
times as fast. A test checks this.

All three models share everything but the social rule: the mission clock, the
stopping limit, the flight layer, perception, $v_0$, and which drones lead.
Their leaders all slow down within $r_{\mathrm{arr}} = 3$\,m of C, by
$s_i = \min(1, |\mathbf g_i|/r_{\mathrm{arr}})$, so they stop there.

\paragraph{Reynolds.} The rules of \citet{reynolds1987} give an acceleration:
\begin{equation}
  \mathbf a_i =
    k_s \!\!\sum_{j:\,d_{ij} < r_s}\!\! \Bigl(\frac{r_s}{d_{ij}} - 1\Bigr)(-\hat{\mathbf d}_{ij})
  + k_a \bigl(\overline{\mathbf v}_j - \mathbf v_i\bigr)
  + k_c\, \overline{\mathbf d}_{ij}
  \;\bigl[\, + k_g\,(v_0 s_i \hat{\mathbf g}_i - \mathbf v_i) \text{ for a leader}\bigr],
\end{equation}
separation, alignment with the neighbours' mean velocity, cohesion toward
their centre, and, for leaders, Reynolds' seek and arrive. The wanted
velocity $\mathbf v_d = \mathbf v_i + \tau \mathbf a_i$ makes the velocity
layer fly exactly $\mathbf a_i$. Nothing propels a follower: it moves only as
its neighbours pull it. Parameters: $r_s = 1$\,m, $k_s = 4$\,m/s$^2$,
$k_a = 4$\,s$^{-1}$, $k_c = 1.2$\,s$^{-2}$, $k_g = 4$\,s$^{-1}$.

\paragraph{Couzin.} In the zone model of \citet{couzin2002}, as used by
\citet{couzin2005}, each drone keeps a direction $\hat{\mathbf u}_i$ as its
own state and flies it at $v_0$. If any neighbour is inside the zone of
repulsion $\alpha$, avoiding them is all that counts; otherwise attraction and
alignment weigh each neighbour the same:
\begin{equation}
  \mathbf u_i \leftarrow
  \begin{cases}
    -\sum_{j:\,d_{ij} < \alpha} \hat{\mathbf d}_{ij} & \text{if any } d_{ij} < \alpha,\\[2pt]
    \sum_{j} \hat{\mathbf d}_{ij} + \sum_{j \text{ moving}} \hat{\mathbf v}_j
      + \hat{\mathbf u}_i & \text{otherwise.}
  \end{cases}
\end{equation}
An informed individual adds the goal with weight $\omega$,
$\hat{\mathbf u}_i \leftarrow \mathrm{unit}(\hat{\mathbf u}_i + \omega \hat{\mathbf g}_i)$,
and the wanted velocity is $v_0 s_i \hat{\mathbf u}_i$. Parameters:
$\alpha = 0.4$\,m, $\omega = 1$; the zones of orientation and attraction
merge into the perception radius. Two departures from the papers, both
deliberate: a follower has no direction until a neighbour moves (Couzin's
agents start with random ones, which would scatter the swarm at the start),
and the turning-rate limit and noise are left to the real drone.

\paragraph{Hybrid.} The model of \citet{vasarhelyi2018}, designed for real
drones, with Reynolds' cohesion and Couzin's informed leaders. Its terms are
velocities:
\begin{equation}
  \mathbf v_d = v_0 \hat{\mathbf h}_i
  + p_r \!\!\sum_{j:\,d_{ij} < r_0}\!\! (r_0 - d_{ij})(-\hat{\mathbf d}_{ij})
  + C_f \sum_j \bigl(|\mathbf v_{ij}| - D_{ij}\bigr)^{+}\, \hat{\mathbf v}_{ij}
  + p_c\, \overline{\mathbf d}_{ij},
\end{equation}
cruise, linear repulsion, friction and cohesion, with
$\mathbf v_{ij} = \mathbf v_j - \mathbf v_i$. Friction is Vásárhelyi's
velocity alignment: two neighbours may differ in velocity by
$D_{ij} = \max\bigl(v_s, D(d_{ij} - r_f)\bigr)$, and only the excess is
corrected. The braking curve
\begin{equation}
  D(x) =
  \begin{cases}
    0 & x \le 0,\\
    p\,x & 0 < p\,x < a/p,\\
    \sqrt{2ax - a^2/p^2} & \text{otherwise}
  \end{cases}
\end{equation}
is the highest speed from which a drone can still stop within $x$, at
deceleration $a$, gentler over the last stretch. Alignment is thus strong up
close and slack far away. The cruise heading $\hat{\mathbf h}_i$ is the mean
direction of the neighbours seen moving, the drone's own included, Couzin's
zone of orientation; a leader uses the same mean without itself, mixed with
the goal, $\mathrm{unit}(\hat{\mathbf h}_i + w\hat{\mathbf g}_i)$, and cruises at
$v_0 s_i$. Parameters, scaled from the authors' simulator to our speed and
spacing and then tuned: $r_0 = 1$\,m, $p_r = 1$\,s$^{-1}$, $C_f = 0.2$,
$v_s = 0.12\,v_0$, $r_f = 2$\,m, $p = 1$\,s$^{-1}$, $a = 1.6$\,m/s$^2$,
$p_c = 0.4$\,s$^{-1}$, $w = 1$.

\paragraph{How the hybrid got its shape.} The hybrid is the model the project
grew on \texttt{l\_hall}, and each of its departures from Vásárhelyi's model
answers a failure seen there:
\begin{itemize}
  \item \emph{Friction gain 0.2, not 1.} Friction sums over about ten
        neighbours, so at gain 1 a drone tries to correct a velocity
        difference ten times faster than it can tilt; a resting flock
        jiggled at 0.2\,m/s. At 0.2 it settles.
  \item \emph{A stopping limit instead of walls that push.} Vásárhelyi keeps
        drones off walls with virtual agents that push them away. A pillar
        forest then looked like a solid wall, and the push, faster than the
        cruise, bounced the flock back. The stopping limit below replaced it.
  \item \emph{The heading is alignment.} At first each drone cruised along
        its own last heading. Eighteen private headings outvoted two leaders:
        drones that slid into a corner stayed there, and one leader flew away
        from C. Couzin's result relies on followers having no direction of
        their own beyond their neighbours', and his informed individual mixes
        the goal with the social direction, not with its own momentum. With
        the heading taken from the neighbours, the flock wakes from its
        leaders outward.
\end{itemize}

\paragraph{The stopping limit.} Every model's wanted velocity ends in the same
limit: toward each neighbour, beyond that neighbour's own speed away, and
toward each wall, a drone may close only as fast as it could still stop
before a margin $m = 0.3$\,m:
\begin{equation}
  \mathbf v_d \cdot \hat{\mathbf d}_{ij} \le
    \mathbf v_j \cdot \hat{\mathbf d}_{ij} + D(d_{ij} - m),
  \qquad
  \mathbf v_d \cdot \hat{\mathbf w} \le D(d_w - m).
\end{equation}
Each inequality is a half-plane of allowed velocities; a few rounds of
projecting onto them in turn find a velocity inside all of them. The walls
get the last rounds, so where the limits conflict the wall's is the one kept:
a drone caught between a wall and a closing neighbour may end up too close to
the neighbour, but never in the wall. The limit only removes
speed, so nothing bounces, and it bounds whatever a model asks for. On the
real dynamics it kept drones about 0.3\,m from walls; between drones, in a
crowd, it did not always hold (\cref{sec:results}).

\paragraph{Tuning.} The hybrid was tuned over the whole of its development on
\texttt{l\_hall}. To compare it fairly, Couzin and Reynolds were then tuned
on the same map under one protocol, fixed before any run: two settings each,
chosen from how the model had failed, 15 runs each (a grid on one seed, the
best three again on a second), on seeds never used by the experiments, scored
by the drones in C at the end. Couzin's best was $\omega = 1$,
$\alpha = 0.4$\,m: a wider zone of repulsion, which overrides everything
else, held the drones apart too long at the gap. Reynolds' best, goal 4 and
alignment 4, brought 1 drone of 20 to C: its leaders stop at the L's inner
corner, where steering straight at C presses them into a wall, a local
minimum no weight removes. A later check with the plan on as well left
$\omega$ at 1, but showed that a Couzin leader off the walls can leave its
flock behind for good, since nothing holds it to the flock's pace. That is a
property of the model as published, kept and measured
(\cref{sec:results}).

\subsection{The leaders' planning layer}
\label{sec:planner}

A leader that steers straight at C fails wherever that line crosses a wall:
Reynolds' leaders stopped at the L's inner corner, and on \texttt{dead\_end}
every model's flock ends in the cup. The leaders therefore get a deliberative
layer on top of the reactive model, the hybrid architecture of
\citet{arkin1998}: a leader plans a path on the known map and follows it,
while the flocking model stays as it is. Followers are untouched: they don't
know a map exists, so the leadership question stays the same, and only what a
leader knows changes.

\paragraph{What a leader plans around.} The arena's outline and every obstacle
at least 1\,m across. Smaller obstacles, such as \texttt{l\_hall}'s pillars,
are left to the stopping limit, since a flock flows through a forest that a
planner would have to treat as a wall.

\paragraph{Path finding.} No single clearance works: a large one erases the
passages drones fly through, a small one makes the path hug every corner and
drag the flock into it. So the clearance is a cost, not a constraint: A*
\citep{hart1968} on a grid of 0.25\,m cells, 8-connected, where cells closer
than the drones' margin (0.3\,m) to an obstacle are forbidden and a step at
distance $d < c = 2$\,m from the nearest one costs
\begin{equation}
  1 + w \left(\frac{c - d}{c - m}\right)^{2}, \qquad w = 4,
\end{equation}
times its length. In a hall 8\,m wide the centre line is free, a 4\,m gap is
crossed through its middle, and a narrower passage is still allowed, at a
price (\cref{fig:maps}). The path is then pruned to its corners: a shortcut
between two waypoints is kept only if it never comes closer to a wall than the
planned path between them did, so pruning cannot undo the clearance. Turned
down: a visibility graph (it needs the one clearance we can't choose),
Voronoi roadmaps (maximal clearance, long detours in open rooms) and sampling
planners such as RRT (built for many dimensions, and random).

\paragraph{Path following.} The leader aims at the point 2\,m ahead along its
path from wherever the flock has pushed it, pure pursuit \citep{coulter1992};
its progress along the path only moves forward, so a path that doubles back
is not cut short. The route hands the model $\mathbf g_i$: the direction to
that point, with the length of the path still to go, so the models' own
arrive term still stops the leader at C and no model had to change.

\paragraph{Planning again.} A leader first planned once, at go time, from
where it hovered. Replays of the failures on \texttt{dead\_end} showed why
that was not enough: the flock had pushed a leader into the cup, its lookahead
point lay on the other side of the cup's arm, and pursuing it pressed the
leader into the wall. Now, whenever the lookahead point is hidden behind a
kept obstacle, the leader plans again from where it is, at most once a
second. The line-of-sight test is exact: a segment clipped against each box
\citep{liang1984} and its distance to each pillar's centre. On every run that
succeeded with a single plan, replanning never triggered: those runs ended at
the same time to the millisecond. It changed only the failures.

\paragraph{Each leader for itself.} Every leader plans its own path from its
own position, with no exchange with the others: a leader that deliberates
with what it knows. Where an obstacle can be passed on either side, two
leaders may choose different sides. This was expected to split the flock at
worst, with each part keeping a leader; \cref{sec:results} shows that for one
model it does something else.

\todo{Video: dead\_end with the plan off (the flock fills the cup) and on,
with a replan.}
